% Basic Notations
\DeclareMathOperator{\st}{\text{s.t.}}

% Set Theory
\DeclareMathOperator{\card}{card}                      % Cardinality

% General Algebra
\DeclareMathOperator{\N}{\mathbb{N}}                    % Set of Natural Numbers
\DeclareMathOperator{\Z}{\mathbb{Z}}                    % Set of Integers
\DeclareMathOperator{\Q}{\mathbb{Q}}                    % Set of Rational Numbers
\DeclareMathOperator{\R}{\mathbb{R}}                    % Set of Real Numbers
\DeclareMathOperator{\C}{\mathbb{C}}                    % Set of Complex Numbers
\DeclareMathOperator{\F}{\mathbb{F}}                    % Number Field (F)
\DeclareMathOperator{\0}{\mathbf{0}}
\DeclareMathOperator{\spn}{span}
\DeclareMathOperator{\Id}{Id}                           % Identity

% Topology
\DeclareMathOperator{\B}{\mathcal{B}}                   % Topological Basis
\DeclareMathOperator{\interior}{int}                     % Interior
\DeclareMathOperator{\dist}{dist}                        % Distance
\DeclareMathOperator{\diam}{diam}                        % Diameter

% Real Analysis
\newcommand{\Borel}{\hyperref[xmp:borel]{\mathcal{B}}}                  % Borel Sigma-Algebra
\newcommand{\Lebout}{\hyperref[dfn:lebesgue-outer-measure]{\ensuremath{m^*}}}              % Lebesgue Outer Measure
\newcommand{\Leb}{\hyperref[dfn:lebesgue-measure]{\ensuremath{m}}}                   % Lebesgue Measure
\newcommand{\Measurable}{\hyperref[dfn:measurable-set]{\mathcal{M}}}                 % Measurable Set
\newcommand{\Char}[1]{\hyperref[dfn:characteristic-function]{\chi}_{#1}}       % Characteristic Function
\newcommand{\Vitali}[1]{\hyperref[dfn:vitali-cover]{\mathcal{V}}(#1)}       % Characteristic Function
\newcommand{\cae}{\hyperref[dfn:ae]{\,\text{a.e.\ \checkmark}}}                  % almost everywhere
\newcommand{\mlim}{\hyperref[dfn:convergence-in-measure]{\xrightarrow{m}}}              % convergence in measure
\newcommand{\aelim}{\hyperref[xmp:ae-convergence]{\xrightarrow{\text{a.e.}}}}              % convergence almost everywhere
\newcommand{\dini}{\hyperref[dfn:dini-derivative]{\mathrm{D}}}              % Dini derivative
\newcommand{\udif}{\hyperref[dfn:upper-lower-derivatives]{\overline{\mathrm{D}}}}              % upper derivative
\newcommand{\ldif}{\hyperref[dfn:upper-lower-derivatives]{\underline{\mathrm{D}}}}             % lower derivative
\newcommand{\Varia}{\hyperref[dfn:total-variation]{\mathrm{V}}}             % Total Variation
\DeclareMathOperator*{\esssup}{\mathrm{ess\,sup}}              % Essential Supremum

% Hyperrefs
\newcommand{\MeasurableSet}{\hyperref[dfn:measurable-set]{Lebesgue 可测集}}                 % Measurable Set
\newcommand{\Fsigma}{\hyperref[xmp:f-sigma-g-delta]{F-sigma 集合}}              % F-sigma Set Family
\newcommand{\Gdelta}{\hyperref[xmp:f-sigma-g-delta]{G-delta 集合}}              % G-delta Set Family
\newcommand{\Lebesguem}{\hyperref[dfn:lebesgue-measure]{Lebesgue 测度}}                   % Lebesgue Measure
\newcommand{\MeasurableFunc}{\hyperref[dfn:measurable-function]{Lebesgue 可测函数}}           % Measurable Function
\newcommand{\SimpleFunc}{\hyperref[dfn:simple-function]{简单函数}}                     % Simple Function
\newcommand{\AlmostEverywhere}{\hyperref[dfn:ae]{几乎处处}}                 % almost everywhere
\newcommand{\IntegrableFunc}{\hyperref[thm:integrable-iff-measurable]{可积函数}}
\newcommand{\BVFunc}{\hyperref[dfn:bounded-variation-function]{有界变差函数}}                 % Bounded Variation Function
\newcommand{\ACFunc}{\hyperref[dfn:absolutely-continuous-function]{绝对连续函数}}                 % Absolutely Continuous Function
